
Interactive narratives are fundamental to educational games, dialogue systems, and generative entertainment, yet contemporary dialogue systems excel at fluency while lacking explicit narrative goals. When users ask unexpected questions or deviate from the intended story arc, language models typically drift off-story, losing narrative coherence and failing to reach intended climaxes within bounded time windows.

We address this limitation through an integrated system combining three mechanisms: (1) beat-based story representation that decomposes narratives into discrete plot points; (2) director mode steering that classifies each user input into a decision type and guides the character back to the story when needed; and (3) a temporal monitor that adapts story pacing under time constraints.

LLM-as-Judge and controlled user studies show that director mode improves narrative coherence and resistance to user derailment: in the user study, participants rated coherence about two points higher on a five-point scale (3.93 vs.\ 1.71, $n$=14). Temporal monitoring improves perceived pacing under a fixed time budget: in the five-minute user-study scenario, 10 of 14 time-constrained rounds were rated well-paced, compared with 1 of 14 without adaptive pacing. Quality drops at budgets shorter than about five minutes, and later beats are compressed or skipped. Our findings suggest that explicit narrative structure combined with real-time steering preserves story coherence and the character's voice, at some cost to conversational flexibility.



\clearpage
{\KOMAoptions{open=any}\chapter*{Task Description}}
\begin{itemize}
\item Familiarize yourself with controllable storytelling for interactive AI characters and our internal AI 
character pipeline. 
\item  Implement a method that constrains the storytelling in a controllable way. For example, constraining 
story generation based on user inputs, the current beat within the story arc and the remaining time to tell the story. 
\item  Thoroughly evaluate the method on several storytelling scenarios and run a user study for human 
evaluation.
\end{itemize}